% Generated deterministically from the audited Markdown source.
% Responsible author: Carptopus.
\documentclass[11pt]{article}
\usepackage[a4paper,margin=27mm]{geometry}
\usepackage[T1]{fontenc}
\usepackage[utf8]{inputenc}
\usepackage{lmodern}
\usepackage{microtype}
\usepackage{amsmath,amssymb,amsthm,mathtools}
\usepackage{needspace}
\usepackage{xcolor}
\usepackage[colorlinks=true,linkcolor=blue!55!black,citecolor=blue!55!black,urlcolor=blue!65!black]{hyperref}
\usepackage{fancyhdr}

\newtheorem{theorem}{Theorem}[section]
\newtheorem{lemma}[theorem]{Lemma}
\newtheorem{corollary}[theorem]{Corollary}
\newtheorem{claim}[theorem]{Claim}
\theoremstyle{definition}
\newtheorem{definition}[theorem]{Definition}
\theoremstyle{remark}
\newtheorem{remark}[theorem]{Remark}
\numberwithin{equation}{section}

\pagestyle{fancy}
\fancyhf{}
\fancyhead[L]{Carptopus}
\fancyhead[R]{Third symbolic powers of co-chordal edge ideals}
\fancyfoot[C]{\thepage}
\setlength{\headheight}{14pt}
\setlength{\parindent}{0pt}
\setlength{\parskip}{0.42em}
\setlength{\emergencystretch}{4em}

\title{The third symbolic power of a co-chordal edge ideal\\
is componentwise linear}
\author{Carptopus\\\href{mailto:carptopus@163.com}{carptopus@163.com}}
\date{25 August 2026}

\hypersetup{
  pdftitle={The third symbolic power of a co-chordal edge ideal is componentwise linear},
  pdfauthor={Carptopus},
  pdfsubject={Componentwise linearity of third symbolic powers of co-chordal edge ideals},
  pdfkeywords={edge ideal, symbolic power, componentwise linear, co-chordal graph, chordal graph, linear resolution}
}

\begin{document}
\maketitle

\begin{center}
\fcolorbox{black!35}{black!4}{\parbox{0.88\textwidth}{\centering\small
Public Beta manuscript, version 0.1-beta. The mathematical claims have passed
project-internal zero-trust proof and manuscript audits; external review is pending.}}
\end{center}

\begin{abstract}


Let $G$ be a finite co-chordal graph and let $I(G)$ be its edge ideal over an
arbitrary field. Recent work of Ahmed and Namiq determines the regularity of
every symbolic power of $I(G)$ and shows that symbolic powers need not be
componentwise linear from the fourth power onward. For the third symbolic
power, their general degree bounds leave exactly one component undecided, the
component generated in degree five. We prove that this component always has a
five-linear resolution. Consequently, $I(G)^{(3)}$ is componentwise linear for
every co-chordal graph. The proof establishes a more general linear-resolution
theorem for a family of fixed-degree monomial ideals defined by capacity-two
constraints on the maximal cliques of a chordal graph and capacity-one
constraints on a bounded collection of marked cliques. The argument combines
simplicial-vertex elimination, polarization, and a two-level induction.

\end{abstract}

\section{Introduction}


Let $K$ be a field and let

\nopagebreak[4]
\[
S=K[x_1,\ldots,x_n].
\]

For a finite simple graph $G$ on $[n]$, its edge ideal is

\nopagebreak[4]
\[
I(G)=(x_i x_j:\{i,j\}\in E(G))\subseteq S.
\]

A homogeneous ideal $J$ is \emph{componentwise linear} if, for every $d$, the ideal
$J_{\langle d\rangle}$ generated by the degree-$d$ elements of $J$ has a
$d$-linear resolution. By Fröberg's theorem \cite{froberg1990}, $I(G)$ has a two-linear
resolution exactly when the complement $G^c$ is chordal. Such graphs $G$ are
called co-chordal.

Ficarra, Moradi, and Römer \cite{ficarra2024} conjectured that every symbolic power of the edge
ideal of a co-chordal graph is componentwise linear. They proved the assertion
for the second symbolic power of every co-chordal graph and for all symbolic
powers in several subclasses. Ahmed and Namiq \cite{ahmed2026} subsequently proved

\nopagebreak[4]
\[
\operatorname{reg} I(G)^{(s)}=2s
\]

for every co-chordal $G$ with an edge and every $s\geq1$. They also disproved
the componentwise-linearity conjecture: for every $s\geq4$, complements of
$s$-sun graphs give a nonlinear component in degree $s+2$.

The third symbolic power is the first power not covered by the general positive
result of \cite{ficarra2024}, but it lies just below the range of the counterexamples in \cite{ahmed2026}.
More precisely, \cite[Theorem 1.2]{ahmed2026} shows for $J=I(G)^{(3)}$ that

\begin{enumerate}
\item $J_{\langle d\rangle}=0$ for $d\leq3$;
\item $J_{\langle4\rangle}$ is zero or has a four-linear resolution;
\item $J_{\langle d\rangle}$ has a $d$-linear resolution for every $d\geq6$.
\end{enumerate}

Thus only $J_{\langle5\rangle}$ remains. Our main result closes this gap.

\Needspace{0.25\textheight}
\begin{theorem}

Let $G$ be a finite co-chordal graph. Over every field,
$I(G)^{(3)}$ is componentwise linear.

\end{theorem}


The proof is independent of the characteristic. Its central statement is a
linear-resolution theorem for monomial ideals attached to a chordal graph
together with at most three marked cliques. The marked cliques record the
capacity drop created when a simplicial vertex is split off. This additional
state makes the elimination recursion close.

The result should be read with two boundaries in mind. First, it proves the
third-power case of the conjecture from \cite{ficarra2024}, not the original all-powers
conjecture, which is false for $s\geq4$ by \cite{ahmed2026}. Second, the proof uses the
special capacity-two structure of the degree-five component. It does not by
itself prove linearity of the highest undecided component in every symbolic
power.

\section{Clique-capacity ideals}


We write $\mathbb N=\{0,1,2,\ldots\}$. If
$a=(a_1,\ldots,a_n)\in\mathbb N^n$, then

\nopagebreak[4]
\[
|a|=\sum_{i=1}^n a_i,
\qquad
a(T)=\sum_{i\in T}a_i
\]

for $T\subseteq[n]$. Let $\mathcal Q(H)$ denote the set of maximal cliques of
a graph $H$.

\Needspace{0.25\textheight}
\begin{definition}

Let $H$ be a chordal graph, and let $\mathcal F$ be a finite collection of
nonempty cliques of $H$. For $2\leq d\leq5$, define

\nopagebreak[4]
\[
M_d(H,\mathcal F)
=
\bigl(x^a:
|a|=d,
\ a(Q)\leq2\text{ for every }Q\in\mathcal Q(H),
\ a(T)\leq1\text{ for every }T\in\mathcal F
\bigr).
\tag{2.1}
\]

Repeated marked cliques may be removed. We also remove empty marked sets after
vertex deletion. The zero ideal is regarded as having a $d$-linear resolution
for every $d$; equivalently, we use $\operatorname{reg}(0)=-\infty$.

\end{definition}


The auxiliary result is the following.

\Needspace{0.25\textheight}
\begin{theorem}[marked-clique capacity theorem]

Let $H$ be a chordal graph, let $2\leq d\leq5$, and let $\mathcal F$ be a
collection of nonempty cliques of $H$. If

\nopagebreak[4]
\[
|\mathcal F|\leq5-d,
\tag{2.2}
\]

then $M_d(H,\mathcal F)$ has a $d$-linear resolution over every field.

\end{theorem}


We first prove the quadratic base case.

\Needspace{0.25\textheight}
\begin{lemma}

If $H$ is any graph and $\mathcal F$ consists of at most three cliques of
$H$, then $M_2(H,\mathcal F)$ has a two-linear resolution.

\end{lemma}


\begin{proof}

The maximal-clique constraints in (2.1) are automatic in degree two. Put

\nopagebreak[4]
\[
U=\bigcup_{T\in\mathcal F}T.
\]

The generators of $M_2(H,\mathcal F)$ are

\begin{enumerate}
\item $x_i x_j$ for $i\neq j$ such that no marked clique contains both $i$ and
   $j$;
\item $x_i^2$ for $i\notin U$.
\end{enumerate}

Polarize this ideal. For every $i\notin U$, write $i'$ for the new vertex in
the polarized generator $x_i x_{i'}$. The polarization is the edge ideal of a
simple graph $E$. We prove that the complement $R=E^c$ is chordal.

On the original vertices, $R$ is the union of the complete graphs on the
members of $\mathcal F$. Suppose that $R$ induced a cycle of length at least
four on original vertices. A marked clique can contain at most one edge of
this cycle: if it contained two cycle edges, then it would also contain two
nonconsecutive cycle vertices and hence give a chord. Since there are at most
three marked cliques, they cannot cover every edge of such a cycle. Therefore
the original-vertex subgraph of $R$ is chordal.

The clone vertices form a clique in $R$. Moreover, $i'$ is adjacent in $R$ to
every original vertex except $i$. Since $i\notin U$, the original vertex $i$
is isolated in the original-vertex subgraph of $R$.

Consider an induced cycle $C$ of length at least four containing a clone
$i'$. If $i\notin V(C)$, then $i'$ is adjacent to every other vertex of $C$
and gives a chord. If both $i$ and $i'$ lie on $C$, then $i$ is not adjacent
to $i'$, and the two neighbors of $i$ on $C$ must be other clones. Those two
clones are adjacent, again giving a chord. Hence $R$ has no induced cycle of
length at least four.

By Fröberg's theorem \cite{froberg1990}, the polarized ideal has a two-linear resolution.
Polarization preserves graded Betti numbers \cite[Section 1.6]{herzog2011}, so
$M_2(H,\mathcal F)$ has a two-linear resolution.
\end{proof}


\section{Simplicial-vertex splitting}


We now prove Theorem 2.2. The following elementary observation will be used
implicitly. If $W$ is obtained by deleting vertices from a graph $H$, then
the capacity-two restrictions on all maximal cliques of $W$ are equivalent to
the restrictions inherited from all maximal cliques of $H$: every clique of
$W$ is contained in a maximal clique of $W$, and every maximal clique of $W$
is contained in a maximal clique of $H$.

\begin{proof}[Proof of Theorem 2.2]

We induct first on $d$, starting with Lemma 2.3, and for fixed $d$ we induct
on $|V(H)|$. If $d\geq3$ and $H$ is complete, its unique maximal clique has
capacity at most two, so $M_d(H,\mathcal F)=0$.

Assume that $H$ is not complete. A standard consequence of Dirac's theorem
for chordal graphs \cite{dirac1961} gives two nonadjacent simplicial vertices of $H$. Since
$|\mathcal F|\leq5-d\leq2$, each marked clique contains at most one of these
two vertices. Therefore one of them belongs to at most one marked clique.
Fix such a simplicial vertex $v$, and put

\nopagebreak[4]
\[
Q=N_H[v].
\]

The set $Q$ is the unique maximal clique of $H$ containing $v$.

Polarize $M_d(H,\mathcal F)$ and let $y$ be the first polarized copy of
$x_v$. Partition the degree-$d$ generators according to whether they contain
$y$. We obtain

\nopagebreak[4]
\[
\operatorname{pol} M_d(H,\mathcal F)=yA+B,
\tag{3.1}
\]

where $B$ comes from exponent vectors with $a_v=0$, and $A$ comes from those
with $a_v\geq1$ after division by $y$.

Define

\nopagebreak[4]
\[
\mathcal F_0
=
\{T\setminus\{v\}:T\in\mathcal F,
\ T\setminus\{v\}\neq\varnothing\}.
\]

The inherited maximal-clique restrictions are equivalent to the maximal
clique restrictions of $H-v$, and hence

\nopagebreak[4]
\[
B=\operatorname{pol}M_d(H-v,\mathcal F_0),
\tag{3.2}
\]

after omitting unused variables. By induction on the number of vertices, $B$
has a $d$-linear resolution.

We next identify $A$. There are two cases.

\textbf{Case 1: $v$ belongs to no marked clique.} Removing one copy of $v$ lowers
only the remaining capacity on $Q$, from two to one. Thus, after relabeling
polarized copies,

\nopagebreak[4]
\[
A=\operatorname{pol}M_{d-1}(H,\mathcal F\cup\{Q\}).
\tag{3.3}
\]

The number of marked cliques satisfies

\nopagebreak[4]
\[
|\mathcal F|+1\leq6-d=5-(d-1),
\]

so the induction on $d$ gives a $(d-1)$-linear resolution of $A$.

\textbf{Case 2: $v$ belongs to a unique marked clique $T$.} Since
$a(T)\leq1$ and $a_v\geq1$, necessarily

\nopagebreak[4]
\[
a_v=1,
\qquad
a(T\setminus\{v\})=0.
\]

All remaining variables therefore lie in $H-T$. Put

\nopagebreak[4]
\[
\mathcal F_1
=
\{R\setminus T:R\in\mathcal F,
\ R\neq T,
\ R\setminus T\neq\varnothing\}
\cup
\{Q\setminus T:Q\setminus T\neq\varnothing\}.
\tag{3.4}
\]

Then

\nopagebreak[4]
\[
A=\operatorname{pol}M_{d-1}(H-T,\mathcal F_1).
\tag{3.5}
\]

Indeed, the old maximal-clique restrictions become the capacity-two maximal
clique restrictions of $H-T$; the old marked clique $R\neq T$ becomes
$R\setminus T$; and the unique maximal clique $Q$ containing $v$ contributes
the capacity-one restriction on $Q\setminus T$. Moreover,

\nopagebreak[4]
\[
|\mathcal F_1|
\leq(|\mathcal F|-1)+1
=|\mathcal F|
\leq5-d
<5-(d-1).
\]

Again, induction on $d$ gives a $(d-1)$-linear resolution of $A$.

It remains to establish the compatibility needed for (3.1).

\Needspace{0.25\textheight}
\begin{claim}

$B\subseteq A$.

\end{claim}


\par\medskip\noindent\emph{Proof of Claim 3.1.}\ 

Let $x^b$ be a degree-$d$ generator of $B$. We find a degree-$(d-1)$
generator $x^c$ of the corresponding unpolarized ideal for $A$ such that
$x^c\mid x^b$.

In Case 1, if $b(Q)=2$, decrease by one a positive coordinate indexed by
$Q$; if $b(Q)\leq1$, decrease any positive coordinate. The resulting vector
$c$ satisfies $c(Q)\leq1$, and every other upper bound remains valid.

In Case 2, note that $T\subseteq Q$ because $v$ is simplicial and $T$ is a
clique containing $v$. There are three subcases.

\begin{enumerate}
\item If $b(T)=1$, decrease the unique positive contribution on $T$. Then
   $c(T)=0$ and $c(Q)\leq1$.
\item If $b(T)=0$ and $b(Q)=2$, decrease a positive coordinate on $Q\setminus T$.
   Such a coordinate exists.
\item If $b(T)=0$ and $b(Q)\leq1$, decrease any positive coordinate.
\end{enumerate}

In every subcase, $c$ is allowed by (3.5) and divides $b$. When passing to
polarizations, delete the highest-numbered copy of the decreased variable.
This gives

\nopagebreak[4]
\[
\operatorname{pol}(x^c)\mid\operatorname{pol}(x^b),
\]

so $B\subseteq A$.
\hfill\qedsymbol
\par\medskip


The generators of $B$ do not contain $y$, and hence $(B:y)=B$. By Claim 3.1,

\nopagebreak[4]
\[
yA\cap B
=y\bigl(A\cap(B:y)\bigr)
=yB.
\tag{3.6}
\]

Consequently there is a short exact sequence

\nopagebreak[4]
\[
0\longrightarrow yB
\longrightarrow yA\oplus B
\longrightarrow \operatorname{pol}M_d(H,\mathcal F)
\longrightarrow0.
\tag{3.7}
\]

The induction hypotheses and the convention $\operatorname{reg}(0)=-\infty$
give

\nopagebreak[4]
\[
\operatorname{reg}(yA)\leq d,
\qquad
\operatorname{reg}(B)\leq d,
\qquad
\operatorname{reg}(yB)\leq d+1.
\]

The standard regularity inequality for (3.7) now gives

\nopagebreak[4]
\[
\operatorname{reg}\operatorname{pol}M_d(H,\mathcal F)\leq d.
\]

If this ideal is nonzero, it is generated in degree $d$, so its regularity is
at least $d$. It therefore has a $d$-linear resolution. Depolarization
preserves the graded Betti numbers, completing the induction.
\end{proof}


\section{The third symbolic power}


We apply Theorem 2.2 to the degree-five component. Let $G$ be co-chordal and
put $H=G^c$. If $G$ has no edge, then $I(G)^{(3)}=0$, so there is nothing to
prove. Assume that $G$ has an edge.

The minimal vertex covers of $G$ are the complements of the maximal cliques
of $H$. Thus \cite[Equation (1)]{ahmed2026} gives

\nopagebreak[4]
\[
I(G)^{(3)}
=
\bigcap_{Q\in\mathcal Q(H)}
(x_i:i\notin Q)^3.
\tag{4.1}
\]

For a monomial $x^a$ of total degree five, membership in (4.1) is equivalent
to

\nopagebreak[4]
\[
\sum_{i\notin Q}a_i\geq3
\quad\text{for every }Q\in\mathcal Q(H),
\]

or, equivalently,

\nopagebreak[4]
\[
a(Q)\leq2
\quad\text{for every }Q\in\mathcal Q(H).
\tag{4.2}
\]

It follows that

\nopagebreak[4]
\[
I(G)^{(3)}_{\langle5\rangle}=M_5(H,\varnothing).
\tag{4.3}
\]

By Theorem 2.2, this ideal has a five-linear resolution. Combining (4.3) with
\cite[Theorem 1.2]{ahmed2026} yields Theorem 1.1.
\hfill\qedsymbol


\Needspace{0.25\textheight}
\begin{corollary}

For co-chordal edge ideals, the componentwise-linearity conjecture of \cite{ficarra2024}
holds for integers $1\leq s\leq3$ over every field, whereas \cite{ahmed2026} gives counterexamples for
every $s\geq4$.

\end{corollary}


\begin{proof}

The case $s=1$ is Fröberg's theorem, the case $s=2$ is \cite[Theorem 3.5]{ficarra2024}, and
the case $s=3$ is Theorem 1.1. The counterexamples for every $s\geq4$ are
given by \cite[Proposition 6.4]{ahmed2026}.
\end{proof}


\section{Remarks on the proof boundary}


The marked-clique capacity theorem is designed for the capacity-two system in
(4.2). Splitting a simplicial vertex lowers one capacity from two to one, and
the bound $|\mathcal F|\leq5-d$ ensures that the induction reaches a quadratic
ideal controlled by at most three marked cliques.

For a higher symbolic power, the analogous highest undecided component has
larger clique capacity. Repeated splitting then creates several capacity
levels rather than a single family of capacity-one marked cliques. The present
argument does not establish closure of that larger state space. Thus no claim
about all highest undecided components, or about any fourth or higher symbolic
power, is made here.

An exact finite implementation was used during discovery to compare both
branches of the splitting identity on all noncomplete nonisomorphic chordal
graphs with at most six vertices and all allowed marked-clique families. It
checked 21,023 states without a failure. This computation is not used in the
proof of Theorem 2.2.

\section*{Acknowledgment of AI assistance}


OpenAI Codex was used as an AI-assisted research and writing tool for
literature triage, symbolic experimentation, proof stress-testing, and
manuscript editing. Carptopus selected the claims, checked the retained
arguments and sources, and assumes responsibility for the manuscript.

\begin{thebibliography}{9}
\footnotesize
\raggedright

\bibitem{ahmed2026}
C.~Ahmed and M.~R.~Namiq,
\emph{Regularity of Symbolic Powers of Co-Chordal Edge Ideals},
arXiv:2608.20542v1 (2026).
\href{https://arxiv.org/abs/2608.20542}{arXiv:2608.20542}.

\bibitem{ficarra2024}
A.~Ficarra, S.~Moradi, and T.~R\"omer,
\emph{Componentwise linear symbolic powers of edge ideals and Minh's conjecture},
arXiv:2411.11537v1 (2024).
\href{https://arxiv.org/abs/2411.11537}{arXiv:2411.11537}.

\bibitem{froberg1990}
R.~Fr\"oberg,
\emph{On Stanley--Reisner rings},
in: Topics in Algebra, Part 2, Banach Center Publications \textbf{26} (1990),
57--70.

\bibitem{herzog2011}
J.~Herzog and T.~Hibi,
\emph{Monomial Ideals},
Graduate Texts in Mathematics 260, Springer, 2011.

\bibitem{dirac1961}
G.~A.~Dirac,
\emph{On rigid circuit graphs},
Abhandlungen aus dem Mathematischen Seminar der Universit\"at Hamburg
\textbf{25} (1961), 71--76.

\end{thebibliography}

\end{document}
